\documentclass[a4paper]{article}
\usepackage{amsmath, bm, xcolor, tikz, algpseudocode, algorithm, natbib, mathtools, mdframed}
\usetikzlibrary{arrows, positioning, backgrounds, fit}

% formatting
\usepackage[margin=1in]{geometry}
\setlength{\parskip}{\baselineskip}
\setlength{\parindent}{0pt}

\title{(Auxiliary) particle filters}
\author{TJ McKinley}
\date{}

% custom commands for brevity
\newcommand{\btheta}{\bm{\theta}}
\newcommand{\bx}{\bm{x}}
\newcommand{\bX}{\bm{X}}
\newcommand{\by}{\bm{y}}
\newcommand{\bz}{\bm{z}}
\newcommand{\bA}{\bm{A}}
\newcommand{\bp}{\bm{p}}
\newcommand{\bq}{\bm{q}}
\newcommand{\br}{\bm{r}}
\newcommand{\bY}{\bm{Y}}
\newcommand{\bZ}{\bm{Z}}
\newcommand{\bU}{\bm{U}}
\newcommand{\bD}{\bm{D}}
\newcommand{\bR}{\bm{R}}
\newcommand{\bI}{\bm{I}}
\newcommand{\bP}{\bm{P}}
\newcommand{\bE}{\bm{E}}
\newcommand{\bH}{\bm{H}}
\newcommand{\bS}{\bm{S}}
\newcommand{\gap}[1][1]{\vspace{#1\baselineskip}}
\newcommand{\srm}[1]{\mbox{\scriptsize #1}}

\bibliographystyle{asa}

\begin{document}

\maketitle

\section{Likelihood terms}

Let's first consider the form of the likelihood functions for the simulator and the model discrepancy terms. Assume that we have $N_s$ spatial regions and $N_a$ age-classes, then for $t = 1, \dots, T$ time points, let $\bz_t = \left\{\bz_{tsa}; s = 1, \dots, N_s~\mbox{and}~a = 1, \dots, N_a\right\}$, where $\bz_{tsa} = \left(D^{Z^\prime}_{Htsa}, D^{Z^\prime}_{Itsa}\right)$ be the \textbf{observed} death incidence counts. Similarly,  $\bx_t = \left\{\bx_{tsa}; s = 1, \dots, N_s~\mbox{and}~a = 1, \dots, N_a\right\}$, where \newline
$\bx_{tsa} = \left(E^X_{tsa}, A^X_{tsa}, R^X_{At}, P^X_{tsa}, I^X_{1tsa}, I^X_{2tsa}, R^X_{Itsa}, D^X_{Itsa}, H^X_{tsa}, D^X_{Htsa}, R^X_{Htsa}\right)$ is a vector of counts from the \textbf{simulator} at time $t$; and $\by_t = \left\{\bz_{tsa}; s = 1, \dots, N_s~\mbox{and}~a = 1, \dots, N_a\right\}$, where \newline $\by_{tsa} = \left(E^Y_{tsa}, A^Y_{tsa}, R^Y_{Atsa}, P^Y_{tsa}, I^Y_{1tsa}, I^Y_{2tsa}, R^Y_{Itsa}, D^Y_{Itsa}, H^Y_{tsa}, D^Y_{Htsa}, R^Y_{Htsa}\right)$ is a vector of counts \textbf{adjusted for model discrepancy} at time $t$. 

We note that some of the likelihood terms below depend on \textbf{incidence} counts (e.g. $D^{X^\prime}_{Htsa}$), but these can be directly derived from the raw counts if required. In the discussion below we use $p\left(\cdot\right)$ to denote an arbitrary probability \textit{mass} function, and $\pi\left(\cdot\right)$ to denote an arbitrary probability \textit{density} function. We use the superscript $T$ to denote \textit{truncated} distributions e.g. $p^T\left(\cdot\right)$ or $\pi^T\left(\cdot\right)$. 

To keep consistent with much of the particle filtering literature, we will use $g\left(\cdot\right)$ to denote a generic observation distributon across all observations, and $f\left(\cdot\right)$ to denote the state transition distributions. In all discussions below we drop dependence on the parameters, $\btheta$, for brevity. Since we have fixed initial conditions, the target likelihood is thus:
\begin{equation}
    \prod_{t = 1}^T g\left(\bz_{t} \mid \by_t, \bx_t\right)f\left(\by_{t}, \bx_t \mid \by_{t - 1}\right) = \prod_{t = 1}^T g\left(\bz_{t} \mid \by_t\right)f\left(\by_{t} \mid \bx_{t}, \by_{t - 1}\right)f\left(\bx_t \mid \by_{t - 1}\right).
\end{equation}
Letting $\bD^{X^\prime}_{Ht} = \left\{D^{X^\prime}_{Htsa}; s = 1, \dots, N_s~\mbox{and}~a = 1, \dots, N_a\right\}$ (and similarly for the other stages), we can then write the simulator distribution as:
\begin{align}
    f\left(\bx_t \mid \by_{t - 1}\right) &= p\left(\bD^{X^\prime}_{Ht}, \bR^{X^\prime}_{Ht} \mid \bH^{Y}_{t - 1}\right)p\left(\bH^{X^\prime}_t, \bI^{X^\prime}_{2t}, \bD^{X^\prime}_{It} \mid \bI^{Y}_{1(t - 1)}\right)p\left(\bR^{X^\prime}_{It} \mid \bI^{Y}_{2(t - 1)}\right)\nonumber\\
    & \qquad \qquad \qquad \times p\left(\bI^{X^\prime}_{1t} \mid \bP^{Y}_{t - 1}\right)p\left(\bP^{X^\prime}_t, \bA^{X^\prime}_t \mid \bE^{Y}_{t - 1}\right)p\left(\bR^{X^\prime}_{At} \mid \bA^{Y}_{t - 1} \right)\nonumber\\
    & \qquad \qquad \qquad \qquad \qquad \qquad \times p\left(\bE^{X^\prime}_t \mid \bS^{Y}_{t - 1}, \bA^{Y}_{t - 1}, \bP^{Y}_{t - 1}, \bI^{Y}_{1(t - 1)}, \bI^{Y}_{2(t - 1)}\right).\label{eq:simlike}
\end{align}
The infection distribution, $p\left(\bE^{X^\prime}_t \mid \bS^{Y}_{t - 1}, \bA^{Y}_{t - 1}, \bP^{Y}_{t - 1}, \bI^{Y}_{1(t - 1)}, \bI^{Y}_{2(t - 1)}\right)$ is a joint distribution over all spatial regions and age-classes, whereas all other transitions are independent over region and age-classes. For example:
\begin{equation}
    p\left(\bD^{X^\prime}_{Ht}, \bR^{X^\prime}_{Ht} \mid \bH^{Y}_{t - 1}\right) = \prod_{s = 1}^{N_s} \prod_{a = 1}^{N_a} p\left(D^{X^\prime}_{Htsa}, R^{X^\prime}_{Htsa} \mid H^{Y}_{(t - 1)sa}\right),
\end{equation}
and similarly for the other transitions.

The conditional model discrepancy distribution is:
\begin{align}
    f\left(\by_t \mid \bx_t, \by_{t - 1}\right) &= p\left(\bD^{Y^\prime}_{Ht} \mid \bD^{X^\prime}_{Ht}, \bH^{Y}_{t - 1}\right)p\left(\bD^{Y^\prime}_{It} \mid \bD^{X^\prime}_{It}, \bI^{Y}_{1(t - 1)}\right)p\left(\bR^{Y^\prime}_{Ht} \mid \bR^{X^\prime}_{Ht}, \bD^{Y^\prime}_{Ht}, \bH^{Y}_{t - 1}\right)\nonumber\\
    & \qquad \times p\left(\bH^Y_t \mid \bH^X_t, \bH^{Y}_{t - 1}, \bD^{Y^\prime}_{Ht}, \bR^{Y^\prime}_{Ht}, \bD^{Y^\prime}_{It}, \bI^{Y}_{1(t - 1)}\right) p\left(\bR^{Y^\prime}_{It} \mid \bR^{X^\prime}_{It}, \bI^{Y}_{2(t - 1)}\right)\nonumber\\
    &\qquad \qquad \times p\left(\bI^Y_{2t} \mid \bI^X_{2t}, \bI^{Y}_{2(t - 1)}, \bR^{Y^\prime}_{It}, \bD^{Y^\prime}_{It}, \bH^{Y^\prime}_t, \bI^{Y}_{1(t - 1)}\right)\nonumber\\
    &\qquad \qquad \qquad \times p\left(\bI^Y_{1t} \mid \bI^X_{1t}, \bI^{Y}_{1(t - 1)}, \bI^{Y^\prime}_{2t}, \bD^{Y^\prime}_{It}, \bH^{Y^\prime}_t, \bP^{Y}_{t - 1}\right)\nonumber\\
    &\qquad \qquad \qquad \qquad \times p\left(\bP^Y_t \mid \bP^X_t, \bP^{Y}_{t - 1}, \bI^{Y^\prime}_{1t}, \bE^{Y}_{t - 1}\right)p\left(\bR^{Y^\prime}_{At} \mid \bR^{X^\prime}_{At}, \bA^{Y}_{t - 1}\right)\nonumber\\
    & \qquad \qquad \qquad \qquad \qquad \times p\left(\bA^Y_t \mid \bA^X_t, \bA^{Y}_{t - 1}, \bR^{Y^\prime}_{At}, \bE^{Y}_{t - 1}\right)\nonumber\\
    & \qquad \qquad \qquad \qquad \qquad \qquad \times p\left(\bE^Y_t \mid \bE^X_t, \bE^{Y}_{t - 1}, \bP^{Y^\prime}_t, \bA^{Y^\prime}_t, \bS^{Y}_{t - 1}\right),\label{eq:mdlike}
\end{align}
and here all the component distributions in (\ref{eq:mdlike}) are independent over space and age given  $\bx_t$ and $\by_{t - 1}$.

Finally, the conditional observation density is:
\begin{align}
    g\left(\bz_t \mid \by_t\right) &= p\left(\bD^{Z^\prime}_{Ht} \mid \bD^{Y^\prime}_{Ht}\right)p\left(\bD^{Z^\prime}_{It} \mid \bD^{Y^\prime}_{It}\right),\label{eq:obsdens}
\end{align}
and all the component distributions in (\ref{eq:obsdens}) are independent over space and age given $\by_t$.

Details on the functional forms for the model discrepancy terms are given in Section~\ref{sec:MD}.

\section{Bootstrap particle filter}

To run a standard bootstrap particle filter (BPF) we can thus sample (with a slight abuse of notation):
\begin{enumerate}
    \item Particle $k$ with probability $\pi^{1:K}_{t - 1}$.
    \item $\bx_t \sim f\left(\bx_t \mid \by^k_{t - 1}\right)$;
    \item $\by_t \sim f\left(\by_t \mid \bx_t, \by^k_{t - 1}\right)$,
\end{enumerate}
and then generating particle weight:
\begin{equation}
    w^k_t = \prod_{s = 1}^{N_s} \prod_{a = 1}^{N_a} \prod_{c = 1}^{N_z} g\left(z_{tsac} \mid y^k_{tsac}\right).
\end{equation}

\section{Auxiliary particle filter}

To simplify future discussions, note that in (\ref{eq:simlike}), we can decompose\footnote{Since $p\left(D^{X^\prime}_{Htsa}, R^{X^\prime}_{Htsa} \mid H^{Y}_{(t - 1)sa}\right)$ is multinomial, both the marginal, $p\left(D^{X^\prime}_{Htsa} \mid H^{Y}_{(t - 1)sa}\right)$, and conditional, $p\left(R^{X^\prime}_{Htsa} \mid D^{X^\prime}_{Htsa}, H^{Y}_{(t - 1)sa}\right)$, are binomial. Similarly the marginal, $p\left(D^{X^\prime}_{Itsa} \mid I^{Y}_{(t - 1)sa}\right)$, is binomial; and the conditional, $p\left(H^{X^\prime}_{tsa}, I^{X^\prime}_{2tsa} \mid D^{X^\prime}_{Itsa}, I^{Y}_{(t - 1)sa}\right)$, is multinomial}
\begin{equation}
    p\left(D^{X^\prime}_{Htsa}, R^{X^\prime}_{Htsa} \mid H^{Y}_{(t - 1)sa}\right) = p\left(D^{X^\prime}_{Htsa} \mid H^{Y}_{(t - 1)sa}\right)p\left(R^{X^\prime}_{Htsa} \mid D^{X^\prime}_{Htsa}, H^{Y}_{(t - 1)sa}\right),
\end{equation}
and
\begin{equation}
    p\left(H^{X^\prime}_{tsa}, I^{X^\prime}_{2tsa}, D^{X^\prime}_{Itsa} \mid I^{Y}_{(t - 1)sa}\right) = p\left(D^{X^\prime}_{Itsa} \mid I^{Y}_{(t - 1)sa}\right)p\left(H^{X^\prime}_{tsa}, I^{X^\prime}_{2tsa} \mid D^{X^\prime}_{Itsa}, I^{Y}_{(t - 1)sa}\right).\label{eq:condsampler}
\end{equation}
Hence if we let $\by^z_t = \left\{D^Y_{Htsa}, D^Y_{Itsa}; s = 1, \dots, N_s~\mbox{and}~a = 1,\dots, N_a\right\}$ be the adjusted states corresponding to the observed states $\bz_t$, and similarly for $\bx^z_t$, then based on (\ref{eq:simlike})--(\ref{eq:condsampler}) above we can then write the likelihood as:
\begin{equation}
    \prod_{t = 1}^T g\left(\bz_{t} \mid \by^z_t\right)f\left(\by^{-z}_{t} \mid \by^z_t, \bx^{-z}_{t}, \by_{t - 1}\right)f\left(\by^z_t \mid \bx^z_t, \by_{t - 1}\right)f\left(\bx^{-z}_t \mid \bx^z_t, \by_{t - 1}\right)f\left(\bx^z_t \mid \by_{t - 1}\right)\label{eq:conddecomp},
\end{equation}
where $\by^{-z}_t = \by_t \setminus \by^z_t$ are those states in $\by_t$ that do not include those in $\by^z_t$, and similarly $\bx^{-z}_t = \bx_t \setminus \bx^z_t$. In the subsequent discussion we focus on importance samplers targeting $f\left(\bx^z_t, \by^z_t \mid \by_{t - 1}\right) = f\left(\by^z_t \mid \bx^z_t, \by_{t - 1}\right)f\left(\bx^z_t \mid \by_{t - 1}\right)$.

From the decomposition in (\ref{eq:conddecomp}), we note that at time $t$ we can write the likelihood contribution as:
\begin{align}
    &g\left(\bz_{t} \mid \by^z_t\right)f\left(\by^{-z}_{t} \mid \by^z_t, \bx^{-z}_{t}, \by_{t - 1}\right)f\left(\by^z_t \mid \bx^z_t, \by_{t - 1}\right)f\left(\bx^{-z}_t \mid \bx^z_t, \by_{t - 1}\right)f\left(\bx^z_t \mid \by_{t - 1}\right) \nonumber\\
    & \qquad \qquad = f\left(\by^{-z}_{t} \mid \by^z_t, \bx^{-z}_{t}, \by_{t - 1}\right)f\left(\bx^{-z}_{t} \mid \bx^z_{t}, \by_{t - 1}\right)\nonumber\\
    & \qquad \qquad \qquad \qquad \times \prod_{s = 1}^{N_s} \prod_{a = 1}^{N_a} \prod_{c = 1}^{N_z} g\left(z_{tsac} \mid y^z_{tsac}\right)f\left(y^z_{tsac} \mid x^z_{tsac}, \by_{(t - 1)sac}\right)f\left(x^z_{tsac} \mid \by_{(t - 1)sac}\right),\label{eq:condlikesa}
\end{align}
where $N_z = 2$ observed states: $D_H$ and $D_I$. Sampling from the first two terms in (\ref{eq:condlikesa}) is described in Sections~\ref{sec:condsim} and \ref{sec:condmdsim}, and hence we focus on an importance sampler:
\begin{equation}
    q\left(y^z_{tsac} \mid z_{tsac}, x^z_{tsac}, \by_{(t - 1)sac}\right) = \frac{g\left(z_{tsac} \mid y^z_{tsac}\right)f\left(y^z_{tsac} \mid x^z_{tsac}, \by_{(t - 1)sac}\right)}{\sum_{y = 0}^n g\left(z_{tsac} \mid y\right)f\left(y \mid x^z_{tsac}, \by_{(t - 1)sac}\right)},
\end{equation}
where all spatial regions, age-classes and states can be sampled independently, and $n$ is the maximum value for $y^z_{tsac}$ given $\by_{(t - 1)sac}$ and $x^z_{tsac}$. From the decomposition in (\ref{eq:conddecomp}) it is clear that we can sample (with a slight abuse of notation):
\begin{enumerate}
    \item Particle $k$ with probability $\pi^{1:K}_{t - 1}$.
    \item $\bx^z_t \sim f\left(\bx^z_t \mid \by^k_{t - 1}\right)$;
    \item $\by^z_t \sim q\left(\by^z_t \mid \bx^z_t, \by^k_{t - 1}\right)$;
    \item $\bx^{-z}_t \sim f\left(\bx^{-z}_t \mid \bx^z_t, \by^k_{t - 1}\right)$;
    \item $\by^{-z}_t \sim f\left(\by^{-z}_t \mid \by^z_t, \bx^{-z}_t, \by^k_{t - 1}\right)$.
\end{enumerate}
With this sampler the particle weights are thus:
\begin{equation}
    w^k_t = \prod_{s = 1}^{N_s} \prod_{a = 1}^{N_a} \prod_{c = 1}^{N_z} \sum_{y = 0}^n g\left(z_{tsac} \mid y\right)f\left(y \mid x^z_{tsac}, \by^k_{(t - 1)sac}\right).
\end{equation}

\section{Dealing with particle impoverishment}

We often have complete particle degeneracy at a given time point $t$, and thus the re-sampled particles have no variability. Following CHOPIN / DOUCETJOHANSEN we can introduce variability by introducing a Markov kernel, $K_t(\cdot)$, of invariant distribution $p\left(\bx_{1:t}, \by_{1:t} \mid \bz_{1:t}\right)$. Such a kernel is only ergodic if all $\bx_{1:t}$ and $\by_{1:t}$ are updated, but following XXXXXXX we will update only $\left(\bx_t, \by_t \mid \bx_{1:t}, \by_{1:t}, \bz_{1:t}\right)$ using independence sampler Metropolis-Hastings update steps. From the model decomposition in (\ref{eq:condlikesa}) we can see that for a given particle we can simulate proposals $\bx_t^\prime$ and $\by_t^\prime$ from the model directly, and then accept-reject the proposal with probability:
\begin{equation}
    \alpha = \min\left(1, \frac{g\left(\bz_{t} \mid \by^{z\prime}_{t}\right)}{g\left(\bz_{tsac} \mid \by^{zk}_t\right)}\right).
\end{equation}
We note also from the decomposition in (\ref{eq:condlikesa}) that $f\left(\bx^{-z}_t \mid \bx^z_t, \by^k_{t - 1}\right)$ and $f\left(\by^{-z}_t \mid \by^z_t, \bx^{-z}_t, \by^k_{t - 1}\right)$ cancel out of the accept-reject ratio, and are independent of $\bx^z_t$ and $\by^z_t$ given $\by^k_{t - 1}$. As such, they do not have to be sampled from in order to accept or reject a proposal. This means that multiple update steps can be done quickly. If, at the end of these at least one proposal has been accepted, then we can simulate $\bx^{-z}_t$ and $\by^{-z}_t$ as required.

For the auxiliary particle filter we sample from $f\left(\bx^z_t \mid \by^k_{t - 1}\right)$ and $q\left(\by^z_t \mid \bx^z_t, \by^k_{t - 1}\right)$ and accept-reject with probability:
\begin{equation}
    \alpha = \min\left(1, \frac{\prod_{s = 1}^{N_s} \prod_{a = 1}^{N_a} \prod_{c = 1}^{N_z} \sum_{y = 0}^{n^\prime_{tsac}} g\left(z_{tsac} \mid y\right)f\left(y \mid x^{z\prime}_{tsac}, \by^k_{(t - 1)sac}\right)}{\prod_{s = 1}^{N_s} \prod_{a = 1}^{N_a} \prod_{c = 1}^{N_z} \sum_{y = 0}^{n^k_{tsac}} g\left(z_{tsac} \mid y\right)f\left(y \mid x^{zk}_{tsac}, \by^k_{(t - 1)sac}\right)}\right).
\end{equation}

\section{Appendix}

\subsection{Sampling from conditional simulator}\label{sec:condsim}

\textcolor{red}{Hence in the subsequent discussion we present the case for sampling from a single region, age-class and state, and for clarity we drop the $s$, $a$ and $c$ subscripts.}

It turns out that we can sample from $\pi\left(\bx^{-z}_t \mid \bx^z_t, \by_{t - 1}\right)$ \textit{directly from the simulator}, by adjusting some parameters and inputs slightly.

If we first consider 
\begin{equation}
    \left(D^{X^\prime}_{Ht}, R^{X^\prime}_{Ht} \mid H^{Y}_{t - 1}\right) \sim \mbox{Multinomial}\left(H^{Y}_{t - 1}, \bp\right),
\end{equation}
where $\bp = \left(p_Hp_{HD}, p_H\left(1 - p_{HD}\right), 1 - p_H\right)$. From this we have that
\begin{equation}
    \left(R^{X^\prime}_{Ht} \mid D^{X^\prime}_{Ht}, H^{Y}_{t - 1}\right) \sim \mbox{Binomial}\left(H^{Y}_{t - 1} - D^{X^\prime}_{Ht}, \frac{p_H\left(1 - p_{HD}\right)}{1 - p_Hp_{HD}}\right).\label{eq:dhbin}
\end{equation}
Therefore (\ref{eq:dhbin}) is equivalent to sampling
\begin{equation}
    \left(D^{\ast X^\prime}_{Ht}, R^{X^\prime}_{Ht} \mid H^{Y}_{t - 1}\right) \sim \mbox{Multinomial}\left(H^{Y}_{t - 1} - D^{X^\prime}_{Ht}, \bp^\ast\right),
\end{equation}
where $\bp^\ast = \left(0, \frac{p_H\left(1 - p_{HD}\right)}{1 - p_Hp_{HD}}, \frac{1 - p_H}{1 - p_Hp_{HD}}\right)$. Similarly, if we consider
\begin{equation}
    \left(H^{X^\prime}_t, I^{X^\prime}_{2t}, D^{X^\prime}_{It} \mid I^{Y}_{1(t - 1)}\right) \sim \mbox{Multinomial}\left(I^{Y}_{1(t - 1)}, \bq\right),
\end{equation}
where $\bq = \left(p_{I_1}p_{I_1H}, p_{I_1}\left(1 - p_{I_1H} - p_{I_1D}\right), p_{I_1}p_{I_1D}, 1 - p_{I_1}\right)$, then we can derive the conditional:
\begin{equation}
    \left(H^{X^\prime}_t, I^{X^\prime}_{2t} \mid D^{X^\prime}_{It}, I^{Y}_{1(t - 1)}\right) \sim \mbox{Multinomial}\left(I^{Y}_{1(t - 1)} - D^{X^\prime}_{It}, \bq^\prime\right),\label{eq:dimult}
\end{equation}
where $\bq^\prime = \left(\frac{p_{I_1}p_{I_1H}}{1 - p_{I_1}p_{I_1D}}, \frac{p_{I_1}\left(1 - p_{I_1H} - p_{I_1D}\right)}{1 - p_{I_1}p_{I_1D}}, \frac{1 - p_{I_1}}{1 - p_{I_1}p_{I_1D}}\right)$. Therefore (\ref{eq:dimult}) is equivalent to sampling
\begin{equation}
    \left(H^{X^\prime}_t, I^{X^\prime}_{2t}, D^{\ast X^\prime}_{It} \mid I^{Y}_{1(t - 1)} - D^{X^\prime}_{It}\right) \sim \mbox{Multinomial}\left(I^{Y}_{1(t - 1)} - D^{X^\prime}_{It}, \bq^\ast\right),
\end{equation}
where $\bq^\ast = \left(\frac{p_{I_1}p_{I_1H}}{1 - p_{I_1}p_{I_1D}}, \frac{p_{I_1}\left(1 - p_{I_1H} - p_{I_1D}\right)}{1 - p_{I_1}p_{I_1D}}, 0, \frac{1 - p_{I_1}}{1 - p_{I_1}p_{I_1D}}\right)$. Hence we can use the simulator directly by slightly adjusting the inputs and parameters to sample from the correct conditional distribution.

\subsection{Sampling from conditional model discrepancy distribution}\label{sec:condmdsim}

The conditional model discrepancy distribution is:
\begin{align}
    f\left(\by^{-z}_t \mid \by^z_t, \bx^{-z}_t, \by_{t - 1}\right) &= p\left(\bR^{Y^\prime}_{Ht} \mid \bR^{X^\prime}_{Ht}, \bD^{Y^\prime}_{Ht}, \bH^{Y}_{t - 1}\right)\nonumber\\
    & \qquad \times p\left(\bH^Y_t \mid \bH^X_t, \bH^{Y}_{t - 1}, \bD^{Y^\prime}_{Ht}, \bR^{Y^\prime}_{Ht}, \bD^{Y^\prime}_{It}, \bI^{Y}_{1(t - 1)}\right) p\left(\bR^{Y^\prime}_{It} \mid \bR^{X^\prime}_{It}, \bI^{Y}_{2(t - 1)}\right)\nonumber\\
    &\qquad \qquad \times p\left(\bI^Y_{2t} \mid \bI^X_{2t}, \bI^{Y}_{2(t - 1)}, \bR^{Y^\prime}_{It}, \bD^{Y^\prime}_{It}, \bH^{Y^\prime}_t, \bI^{Y}_{1(t - 1)}\right)\nonumber\\
    &\qquad \qquad \qquad \times p\left(\bI^Y_{1t} \mid \bI^X_{1t}, \bI^{Y}_{1(t - 1)}, \bI^{Y^\prime}_{2t}, \bD^{Y^\prime}_{It}, \bH^{Y^\prime}_t, \bP^{Y}_{t - 1}\right)\nonumber\\
    &\qquad \qquad \qquad \qquad \times p\left(\bP^Y_t \mid \bP^X_t, \bP^{Y}_{t - 1}, \bI^{Y^\prime}_{1t}, \bE^{Y}_{t - 1}\right)p\left(\bR^{Y^\prime}_{At} \mid \bR^{X^\prime}_{At}, \bA^{Y}_{t - 1}\right)\nonumber\\
    & \qquad \qquad \qquad \qquad \qquad \times p\left(\bA^Y_t \mid \bA^X_t, \bA^{Y}_{t - 1}, \bR^{Y^\prime}_{At}, \bE^{Y}_{t - 1}\right)\nonumber\\
    & \qquad \qquad \qquad \qquad \qquad \qquad \times p\left(\bE^Y_t \mid \bE^X_t, \bE^{Y}_{t - 1}, \bP^{Y^\prime}_t, \bA^{Y^\prime}_t, \bS^{Y}_{t - 1}\right),\label{eq:mdlikecond}
\end{align}
and here all the component distributions in (\ref{eq:mdlikecond}) are independent over space and age given $\by^z_t$, $\bx^{-z}_t$ and $\by_{t - 1}$ and thus are straightforward to sample from.

\subsection{Model discrepancy constraints}\label{sec:MD}

\textcolor{red}{Note order is different to non-spatial model as have swapped DI and H}.

For \textbf{absorbing states} ($D_H$, $R_H$, $D_I$, $R_I$ and $R_A$) we place the model discrepancy on the \textit{incidence} (i.e. new cases) rather than the state counts. For all other states, we place the model discrepancy on the numbers of individuals in each state directly. A key aspect is that we must ensure that any model discrepancy added results in states that are \textit{epidemiologically} valid. Hence we can't have more deaths than infections and such like. This means we add constraints onto the model discrepancy process based on the structure of the model.

Because we're focusing on individual states, we will use the notation e.g. $D^X_{Ht}$ to denote the state counts from the \textit{simulator} at time $t$ (relating to $X_t$ above), $D^Y_{Ht}$ is the \textit{adjusted} state after model discrepancy is added, and e.g. $D^{X^\prime}_{Ht}$ is the \textit{incidence}.

Thinking first about $D^Y_{Ht}$ we have that $D^Y_{Ht} = D^Y_{H(t - 1)} + D^{Y^\prime}_{Ht}$, and hence we can derive bounds:
\begin{equation*}
    0 \leq D^{Y^\prime}_{Ht} \leq H^Y_{t-1} \qquad \mbox{since $D^{Y^\prime}_{Ht} \leq H^Y_{t - 1}$}.
\end{equation*}
We let the adjusted incidence $D^{Y^\prime}_{Ht} = D^{X^\prime}_{Ht} + M^\prime_{D_Ht}$, and then place a truncated Skellam distribution on $M^\prime_{D_Ht}$, with constraints:
\begin{equation*}
    -D^{X^\prime}_{Ht} \leq M^\prime_{D_Ht} \leq H^Y_{t - 1} - D^{X^\prime}_{Ht}.
\end{equation*}
Then the adjusted count is given by
\begin{equation*}
    D^Y_{Ht} = D^X_{Ht} + M^\prime_{D_Ht}.
\end{equation*}
This holds since
\begin{eqnarray*}
    D^Y_{Ht} &=& D^Y_{H(t - 1)} + D^{Y^\prime}_{Ht}\\
    &=& D^Y_{H(t - 1)} + D^{X^\prime}_{Ht} + M^\prime_{D_Ht}\\
    &=& D^X_{Ht} + M^\prime_{D_Ht} \qquad \mbox{since}~D^X_{Ht} = D^Y_{H(t - 1)} + D^{X^\prime}_{Ht}.
\end{eqnarray*}
To summarise:
\begin{align*}
\Aboxed{M^\prime_{D_Ht} &\sim \mbox{Skellam}\left(\lambda_1, \lambda_2\right)I\left(-D^{X^\prime}_{Ht}, H^Y_{t - 1} - D^{X^\prime}_{Ht}\right)}\\
    \Aboxed{D^{Y^\prime}_{Ht} &= D^{X^\prime}_{Ht} + M^\prime_{D_Ht}}\\
    \Aboxed{D^Y_{Ht} &= D^X_{Ht} + M^\prime_{D_Ht}}
\end{align*}

Now consider $D^Y_{It}$. Since this is an absorbing state, we place MD on the \textit{incidence} $D^{Y^\prime}_{It}$, and hence we can derive bounds for the MD as:
\begin{align*}
    0 &\leq D^{Y^\prime}_{It} \leq I^Y_{1(t - 1)}\\
    0 &\leq D^{X^\prime}_{It} + M^\prime_{D_It} \leq I^Y_{1(t - 1)}\\ 
    -D^{X^\prime}_{It} &\leq M^\prime_{D_It} \leq I^Y_{1(t - 1)} - D^{X^\prime}_{It}.
\end{align*}
Hence
\begin{align*}
\Aboxed{M^\prime_{D_It} &\sim \mbox{Skellam}\left(\lambda_1, \lambda_2\right)I\left(-D^{X^\prime}_{It}, I^Y_{1(t - 1)} - D^{X^\prime}_{It}\right)}\\
    \Aboxed{D^{Y^\prime}_{It} &= D^{X^\prime}_{It} + M^\prime_{D_It}}\\
    \Aboxed{D^Y_{It} &= D^X_{It} + M^\prime_{D_It}}
\end{align*}

We also know that
\begin{align*}
    &D^{Y^\prime}_{Ht} + R^{Y^\prime}_{Ht} \leq H^Y_{t - 1}\\
    \Rightarrow 0 &\leq R^{Y^\prime}_{Ht} \leq H^Y_{t - 1} - D^{Y^\prime}_{Ht},
\end{align*}
and hence we can derive a constraint for $M^\prime_{R_Ht}$ conditional on the adjusted $D^{Y^\prime}_{Ht}$, such that
\begin{equation*}
    -R^{X^\prime}_{Ht} \leq M^\prime_{R_Ht} \leq H^Y_{t - 1} - D^{Y^\prime}_{Ht} - R^{X^\prime}_{Ht}.
\end{equation*}
Hence we have:
\begin{align*}
\Aboxed{M^\prime_{R_Ht} &\sim \mbox{Skellam}\left(\lambda_1, \lambda_2\right)I\left(-R^{X^\prime}_{Ht}, H^Y_{t - 1} - D^{Y^\prime}_{Ht} - R^{X^\prime}_{Ht}\right)}\\
    \Aboxed{R^{Y^\prime}_{Ht} &= R^{X^\prime}_{Ht} + M^\prime_{R_Ht}}\\
    \Aboxed{R^Y_{Ht} &= R^X_{Ht} + M^\prime_{R_Ht}}
\end{align*}

Now consider $H^Y_t$. We know that
\begin{equation*}
    H^Y_t = H^Y_{t - 1} + H^{Y^\prime}_t - D^{Y^\prime}_{Ht} - R^{Y^\prime}_{Ht},
\end{equation*}
and we now place the model discrepancy on the state and \textit{not} the incidence, since $H^Y_t$ is not an absorbing state. Therefore, 
\begin{equation*}
    H^Y_t = H^X_t + M_{Ht}.
\end{equation*}
Hence
\begin{align*}
    H^Y_t &= H^Y_{t - 1} + H^{Y^\prime}_t - D^{Y^\prime}_{Ht} - R^{Y^\prime}_{Ht}\\
    \Rightarrow H^{Y^\prime}_t &= H^Y_t - H^Y_{t - 1} + D^{Y^\prime}_{Ht} + R^{Y^\prime}_{Ht}\\
    \Rightarrow H^{Y^\prime}_t &= H^X_t + M_{Ht} - H^Y_{t - 1} + D^{Y^\prime}_{Ht} + R^{Y^\prime}_{Ht}.
\end{align*}
Since $0 \leq H^{Y^\prime}_t \leq I^Y_{1(t - 1)} - D^{Y^\prime}_{It}$ we can thus derive bounds:
\begin{align*}
    0 &\leq H^X_t + M_{Ht} - H^Y_{t - 1} + D^{Y^\prime}_{Ht} + R^{Y^\prime}_{Ht} \leq I^Y_{1(t - 1)} - D^{Y^\prime}_{It}\\
    \Rightarrow -H^X_t + H^Y_{t - 1} - D^{Y^\prime}_{Ht} - R^{Y^\prime}_{Ht} &\leq M_{Ht} \leq I^Y_{1(t - 1)}  - D^{Y^\prime}_{It} - H^X_t + H^Y_{t - 1} - D^{Y^\prime}_{Ht} - R^{Y^\prime}_{Ht}.
\end{align*}
Hence
\begin{align*}
\Aboxed{M_{Ht} &\sim \mbox{Skellam}\left(\lambda_1, \lambda_2\right)I\left(-H^X_t + H^Y_{t - 1} - D^{Y^\prime}_{Ht} - R^{Y^\prime}_{Ht}, I^Y_{1(t - 1)} - D^{Y^\prime}_{It} - H^X_t + H^Y_{t - 1} - D^{Y^\prime}_{Ht} - R^{Y^\prime}_{Ht}\right)}\\
    \Aboxed{H^Y_{t} &= H^X_{t} + M_{Ht}}\\
    \Aboxed{H^{Y^\prime}_{t} &= H^X_t + M_{Ht} - H^Y_{t - 1} + D^{Y^\prime}_{Ht} + R^{Y^\prime}_{Ht}}
\end{align*}

We also have that $0 \leq R^{Y^\prime}_{It} \leq I^{Y^\prime}_{2(t - 1)}$, leading to:
\begin{align*}
\Aboxed{M^\prime_{R_It} &\sim \mbox{Skellam}\left(\lambda_1, \lambda_2\right)I\left(-R^{X^\prime}_{It}, I^Y_{2(t - 1)} - R^{X^\prime}_{It}\right)}\\
    \Aboxed{R^{Y^\prime}_{It} &= R^{X^\prime}_{It} + M^\prime_{R_It}}\\
    \Aboxed{R^Y_{It} &= R^X_{It} + M^\prime_{R_It}}
\end{align*}

Next, we know that
\begin{equation*}
    I^Y_{2t} = I^Y_{2(t - 1)} + I^{Y^\prime}_{2t} - R^{Y^\prime}_{It},
\end{equation*}
and we place MD on the \textit{counts}, such that:
\begin{equation*}
    I^Y_{2t} = I^X_{2t} + M_{I_2t}
\end{equation*}
We know that:
\begin{align*}
    I^Y_{2t} &= I^Y_{2(t - 1)} + I^{Y^\prime}_{2t} - R^{Y^\prime}_{It}\\
    \Rightarrow I^{Y^\prime}_{2t} &= I^Y_{2t} - I^Y_{2(t - 1)} + R^{Y^\prime}_{It}\\
    \Rightarrow I^{Y^\prime}_{2t} &= I^X_{2t} + M_{I_2t} - I^Y_{2(t - 1)} + R^{Y^\prime}_{It}
\end{align*}
Since $0 \leq I^{Y^\prime}_{2t} \leq I^Y_{1(t - 1)} - D^{Y^\prime}_{It} - H^{Y^\prime}_t$, we can derive bounds:
\begin{align*}
    0 &\leq I^{Y^\prime}_{2t} \leq I^Y_{1(t - 1)} - D^{Y^\prime}_{It} - H^{Y^\prime}_t\\
    \Rightarrow 0 &\leq I^X_{2t} + M_{I_2t} - I^Y_{2(t - 1)} + R^{Y^\prime}_{It} \leq I^Y_{1(t - 1)} - D^{Y^\prime}_{It} - H^{Y^\prime}_t\\
    \Rightarrow -I^X_{2t} + I^Y_{2(t - 1)} - R^{Y^\prime}_{It} &\leq M_{I_2t} \leq I^Y_{1(t - 1)} - D^{Y^\prime}_{It} - H^{Y^\prime}_t - I^X_{2t} + I^Y_{2(t - 1)} - R^{Y^\prime}_{It}.
\end{align*}
Hence
\begin{align*}
\Aboxed{M_{I_2t} &\sim \mbox{Skellam}\left(\lambda_1, \lambda_2\right)I\left(-I^X_{2t} + I^Y_{2(t - 1)} - R^{Y^\prime}_{It}, I^Y_{1(t - 1)} - D^{Y^\prime}_{It} - H^{Y^\prime}_t - I^X_{2t} + I^Y_{2(t - 1)} - R^{Y^\prime}_{It}\right)}\\
    \Aboxed{I^Y_{2t} &= I^X_{2t} + M_{I_2t}}\\
    \Aboxed{I^{Y^\prime}_{2t} &= I^X_{2t} + M_{I_2t} - I^Y_{2(t - 1)} + R^{Y^\prime}_{It}}
\end{align*}

Then we can construct $I^Y_{1t}$ by considering
\begin{equation*}
    I^Y_{1t} = I^Y_{1(t - 1)} + I^{Y^\prime}_{1t} - I^{Y^\prime}_{2t} - D^{Y^\prime}_{It} - H^{Y^\prime}_t,
\end{equation*}
and we place MD on the \textit{counts}, such that:
\begin{equation*}
    I^Y_{1t} = I^X_{1t} + M_{I_1t}
\end{equation*}
Hence we have:
\begin{align*}
    I^Y_{1t} &= I^Y_{1(t - 1)} + I^{Y^\prime}_{1t} - I^{Y^\prime}_{2t} - D^{Y^\prime}_{It} - H^{Y^\prime}_t\\
    I^{Y^\prime}_{1t} &= I^Y_{1t} - I^Y_{1(t - 1)} + I^{Y^\prime}_{2t} + D^{Y^\prime}_{It} + H^{Y^\prime}_t
\end{align*}
and since $0 \leq I^{Y^\prime}_{1t} \leq P^Y_{t - 1}$ we can derive bounds:
\begin{align*}
    0 &\leq I^{Y^\prime}_{1t} \leq P^Y_{t - 1}\\
    0 &\leq I^Y_{1t} - I^Y_{1(t - 1)} + I^{Y^\prime}_{2t} + D^{Y^\prime}_{It} + H^{Y^\prime}_t \leq P^Y_{t - 1}\\
    0 &\leq I^X_{1t} + M_{I_1t} - I^Y_{1(t - 1)} + I^{Y^\prime}_{2t} + D^{Y^\prime}_{It} + H^{Y^\prime}_t \leq P^Y_{t - 1}\\
    -I^X_{1t} + I^Y_{1(t - 1)} - I^{Y^\prime}_{2t} - D^{Y^\prime}_{It} - H^{Y^\prime}_t &\leq M_{I_1t} \leq P^Y_{t - 1} - I^X_{1t} + I^Y_{1(t - 1)} - I^{Y^\prime}_{2t} - D^{Y^\prime}_{It} - H^{Y^\prime}_t
\end{align*}
Hence
\begin{align*}
\Aboxed{M_{I_1t} &\sim \mbox{Skellam}\left(\lambda_1, \lambda_2\right)I\left(-I^X_{1t} + I^Y_{1(t - 1)} - I^{Y^\prime}_{2t} - D^{Y^\prime}_{It} - H^{Y^\prime}_t, P^Y_{t - 1} - I^X_{1t} + I^Y_{1(t - 1)} - I^{Y^\prime}_{2t} - D^{Y^\prime}_{It} - H^{Y^\prime}_t\right)}\\
    \Aboxed{I^Y_{1t} &= I^X_{1t} + M_{I_1t}}\\
    \Aboxed{I^{Y^\prime}_{1t} &= I^X_{1t} + M_{I_1t} - I^Y_{1(t - 1)} + I^{Y^\prime}_{2t} + D^{Y^\prime}_{It} + H^{Y^\prime}_t}
\end{align*}

For $P^Y_t$ we can consider that
\begin{equation*}
    P^Y_t = P^Y_{t - 1} + P^{Y^\prime}_t - I^{Y^\prime}_{1t}
\end{equation*}
with MD placed on the \textit{counts} as
\begin{equation*}
    P^Y_t = P^X_T + M_{Pt}.
\end{equation*}
Since $0 \leq P^{Y^\prime}_t \leq E^Y_{t - 1}$, we can derive bounds
\begin{align*}
    0 &\leq P^{Y^\prime}_t \leq E^Y_{t - 1}\\
    0 &\leq P^Y_t - P^Y_{t - 1} + I^{Y^\prime}_{1t} \leq E^Y_{t - 1}\\
    0 &\leq P^X_t + M_{Pt} - P^Y_{t - 1} + I^{Y^\prime}_{1t} \leq E^Y_{t - 1}\\
    -P^X_t + P^Y_{t - 1} - I^{Y^\prime}_{1t} &\leq M_{Pt} \leq E^Y_{t - 1} - P^X_t + P^Y_{t - 1} - I^{Y^\prime}_{1t}.
\end{align*}
Hence
\begin{align*}
\Aboxed{M_{Pt} &\sim \mbox{Skellam}\left(\lambda_1, \lambda_2\right)I\left(-P^X_t + P^Y_{t - 1} - I^{Y^\prime}_{1t}, E^Y_{t - 1} - P^X_t + P^Y_{t - 1} - I^{Y^\prime}_{1t}\right)}\\
    \Aboxed{P^Y_t &= P^X_t + M_{Pt}}\\
    \Aboxed{P^{Y^\prime}_t &= P^X_t + M_{Pt} - P^Y_{t - 1} + I^{Y^\prime}_{1t}}
\end{align*}

For $R^Y_{At}$ we can consider:
\begin{equation*}
    R^Y_{At} = R^Y_{A(t - 1)} + R^{Y^\prime}_{At}
\end{equation*}
and we place MD on the \textit{incidence} as
\begin{equation*}
    R^{Y^\prime}_{At} = R^{X^\prime}_{At} + M^\prime_{R_At}.
\end{equation*}
We can derive bounds as:
\begin{align*}
    0 &\leq R^{Y^\prime}_{At} \leq A^Y_{t - 1}\\
    0 &\leq R^{X^\prime}_{At} + M^\prime_{At} \leq A^Y_{t - 1}\\
    -R^{X^\prime}_{At} &\leq  M^\prime_{At} \leq A^Y_{t - 1} - R^{X^\prime}_{At}
\end{align*}
and thus:
\begin{align*}
\Aboxed{M_{R_At} &\sim \mbox{Skellam}\left(\lambda_1, \lambda_2\right)I\left(-R^{X^\prime}_{At}, A^Y_{t - 1} - R^{X^\prime}_{At}\right)}\\
    \Aboxed{R^{Y^\prime}_{At} &= R^{X^\prime}_{At} + M^\prime_{R_At}}\\
    \Aboxed{R^Y_{At} &= R^X_{At} + M^\prime_{R_At}}
\end{align*}

For $A^Y_t$ we can consider that
\begin{equation*}
    A^Y_t = A^Y_{t - 1} + A^{Y^\prime}_t - R^{Y^\prime}_{At}
\end{equation*}
and place MD on the \textit{counts}, such that
\begin{equation*}
    A^Y_t = A^X_t + M_{At}
\end{equation*}
Since $0 \leq A^{Y^\prime}_t \leq E^Y_{t - 1} - P^{Y^\prime}_t$, we can derive bounds:
\begin{align*}
    0 &\leq A^{Y^\prime}_t \leq E^Y_{t - 1} - P^{Y^\prime}_t\\
    0 &\leq A^Y_t - A^Y_{t - 1} + R^{Y^\prime}_{At} \leq E^Y_{t - 1} - P^{Y^\prime}_t\\
    0 &\leq A^X_t + M_{At} - A^Y_{t - 1} + R^{Y^\prime}_{At} \leq E^Y_{t - 1} - P^{Y^\prime}_t\\
    -A^X_t + A^Y_{t - 1} - R^{Y^\prime}_{At} &\leq M_{At} \leq E^Y_{t - 1} - P^{Y^\prime}_t - A^X_t + A^Y_{t - 1} - R^{Y^\prime}_{At}
\end{align*}
Hence
\begin{align*}
\Aboxed{M_{At} &\sim \mbox{Skellam}\left(\lambda_1, \lambda_2\right)I\left(-A^X_t + A^Y_{t - 1} - R^{Y^\prime}_{At}, E^Y_{t - 1} - P^{Y^\prime}_t - A^X_t + A^Y_{t - 1} - R^{Y^\prime}_{At}\right)}\\
    \Aboxed{A^Y_t &= A^X_t + M_{At}}\\
    \Aboxed{A^{Y^\prime}_t &= A^X_t + M_{At} - A^Y_{t - 1} + R^{Y^\prime}_{At}}
\end{align*}

Then, we have that
\begin{equation*}
    E^Y_t = E^Y_{t - 1} + E^{Y^\prime}_t - P^{Y^\prime}_t - A^{Y^\prime}_t
\end{equation*}
and we place MD on the \textit{count}, such that
\begin{equation*}
    E^Y_t = E^X_t + M_{Et}
\end{equation*}
Since $0 \leq E^{Y^\prime}_t \leq S^Y_{t - 1}$ we can derive bounds:
\begin{align*}
    0 &\leq E^{Y^\prime}_t \leq S^Y_{t - 1}\\
    0 &\leq E^Y_t - E^Y_{t - 1} + P^{Y^\prime}_t + A^{Y^\prime}_t \leq S^Y_{t - 1}\\
    0 &\leq E^X_t + M_{Et} - E^Y_{t - 1} + P^{Y^\prime}_t + A^{Y^\prime}_t \leq S^Y_{t - 1}\\
    -E^X_t + E^Y_{t - 1} - P^{Y^\prime}_t - A^{Y^\prime}_t &\leq M_{Et} \leq S^Y_{t - 1} - E^X_t + E^Y_{t - 1} - P^{Y^\prime}_t - A^{Y^\prime}_t 
\end{align*}
and hence
\begin{align*}
\Aboxed{M_{Et} &\sim \mbox{Skellam}\left(\lambda_1, \lambda_2\right)I\left(-E^X_t + E^Y_{t - 1} - P^{Y^\prime}_t - A^{Y^\prime}_t, S^Y_{t - 1} - E^X_t + E^Y_{t - 1} - P^{Y^\prime}_t - A^{Y^\prime}_t\right)}\\
    \Aboxed{E^Y_t &= E^X_t + M_{Et}}\\
    \Aboxed{E^{Y^\prime}_t &= E^X_t + M_{Et} - E^Y_{t - 1} + P^{Y^\prime}_t + A^{Y^\prime}_t}
\end{align*}

Finally we have
\begin{align*}
    \Aboxed{S^Y_t &= S^Y_{t - 1} - E^{Y^\prime}_t}
\end{align*}

\bibliography{/home/tj/Documents/ref}

\end{document}
